\documentclass{situationist-preprint}

\title{Against Modal Realism}
\msid{TSM-08-a}
\msfamily{TSM-08}
\msdiscipline{Time, Space and Modality}
\mstype{Principal manuscript}
\mskeywords{modal realism; possible worlds; ersatzism; ontological economy; recombination; Lewis}

\begin{document}
\maketitle

\begin{abstract}
Do other possible worlds exist concretely? Lewis argued that they do, on the ground that the hypothesis of a plurality of concrete worlds earns its place by its theoretical benefits, as set theory earns its place in mathematics. We accept the form of the argument and reject its conclusion. We go through the work Lewis assigned to concrete worlds: the analysis of modality, of counterfactuals, of content and of properties. For each, we show that an abstract construction does the same work, and that the one advantage claimed specifically for concrete worlds, a reduction of modality to the non-modal, is not delivered, because the principle of recombination that fixes which worlds there are must be restricted by modal notions. Since concrete worlds are causally and spatiotemporally isolated from us, nothing but theoretical utility could warrant belief in them; an ontology that does no work its rivals cannot do is not warranted, and other worlds do not exist concretely.
\end{abstract}

\section{Lewis's argument}

\citet{Lewis1986} held that there are many concrete universes, each a maximal sum of spatiotemporally interrelated things, none related spatiotemporally or causally to any other, and that ours is distinguished only as the one we inhabit. He did not claim that we observe the others or infer them as causes. His argument was cost-benefit: the hypothesis of a plurality of worlds serves as \enquote{a philosophers' paradise}, giving analyses of modality, counterfactuals, content and properties, and its benefits outweigh its ontological cost, as the benefits of set theory outweigh the cost of believing in sets \citep[\S1.1]{Lewis1986}.

We accept the form of the argument. Concrete worlds are, by hypothesis, isolated from us, and so no evidence bears on them except their contribution to theory. We dispute the premise that they make a contribution that the ersatz alternatives do not. The ersatz alternative we use is that of TSM-07, on which worlds are maximal consistent sets of propositions and the actual world is the only concrete one.

\section{The principle}

\begin{principle}[Economy for isolated entities]\label{p:econ}
If entities of a kind are causally and spatiotemporally isolated from us, belief in them is warranted only if a theory that posits them does explanatory work that no theory without them does at lower ontological cost.
\end{principle}

Principle~\ref{p:econ} is weaker than a general razor: for entities knowable only through their theoretical role, that role is the only warrant. Lewis accepted it, since his defense of the plurality of worlds rested entirely on theoretical benefit.

\section{The work, done without concrete worlds}

\paragraph{Modality.} On modal realism, $p$ is possible if it is true at some world. On the ersatz account, $p$ is possible if it belongs to some maximal consistent set of propositions (TSM-07), and, by TSM-10, consistency is logical consistency with the essential truths. The two accounts validate the same modal logic and assign the same truth conditions to modal sentences, once the extension of \enquote{possible} is fixed.

\paragraph{Counterfactuals.} The similarity semantics of \citet{Lewis1973} requires a domain of worlds ordered by similarity to the actual world. Abstract worlds may be ordered by similarity as well as concrete ones: the similarity of two maximal sets of propositions is a matter of which propositions they share and how important these are, which is what the concrete account appeals to.

\paragraph{Content.} Lewis took the content of a belief to be a set of possible worlds, or of centred worlds. Sets of ersatz worlds serve equally. Where content must be finer than sets of possible worlds, as TSM-09 argues it must, both accounts need impossible worlds, and only the ersatz account can supply them, since a concrete impossible world would be a concrete thing at which a contradiction is true.

\paragraph{Properties.} Lewis identified a property with the set of its possible instances; the ersatzist models it as a function from ersatz worlds to extensions.

\section{The reduction}

There remains the benefit Lewis prized most: modal realism analyses modality without modal primitives, while the ersatzist must say which sets of propositions are consistent and so uses a modal notion \citep[\S3.2]{Lewis1986}.

\begin{proposition}\label{p:recomb}
Modal realism yields a non-modal analysis of modality only if it can state, without modal notions, which worlds there are; and it cannot.
\end{proposition}

\begin{proof}
Modal realism analyses \enquote{possibly $p$} as \enquote{at some world, $p$}. The analysis is correct only if there are enough worlds: one for every way things could be. The principle that says how many there are is Lewis's principle of recombination: anything can coexist with anything, and fail to coexist with anything, \enquote{size and shape permitting} \citep[\S1.8]{Lewis1986}. The proviso is required, since without it recombination generates more duplicates of things than any spacetime can contain, and the proviso says that the recombination must fit into some possible spacetime. Its content is modal. \citet{DiversMelia2002} argued, more generally, that no non-modal statement of the extent of logical space secures that there is a world for every possibility, so that the realist's analysis is either incomplete or modal. The analysis therefore uses a modal notion in fixing its domain.
\end{proof}

\begin{theorem}\label{t:main}
The ontology of concrete worlds does no explanatory work that the ersatz account cannot do. Hence, by Principle~\ref{p:econ}, belief in concrete worlds other than the actual one is not warranted.
\end{theorem}

\begin{proof}
The work Lewis assigned to worlds is the analysis of modality, of counterfactuals, of content and of properties. The ersatz account does each of the last three, and the first except for the reduction of modality to the non-modal, by the preceding section. By Proposition~\ref{p:recomb}, modal realism does not achieve that reduction either. The ersatz account posits no concrete entity other than the cosmos and only abstract entities, propositions and their sets, that every theory of thought requires; its cost is therefore lower. Concrete worlds are isolated from us by hypothesis. Principle~\ref{p:econ} then applies.
\end{proof}

We conclude that other possible worlds do not exist concretely: there is no reason to believe in them, and a reason, economy, to disbelieve. The ersatz account, moreover, alone supplies the impossible worlds of TSM-09.

\section{Objections}

\begin{objection}[The ersatzist's primitive]
Even if modal realism uses a modal notion in its principle of recombination, it uses one fewer than the ersatzist, who needs both consistency and a relation of representation.
\end{objection}

\begin{reply}
The ersatzist's representation is the representational property of propositions, which every theory needs to explain belief and assertion; Lewis supplies it by sets of worlds, a construction that presupposes modal realism and cannot be cited in its favor. The ersatzist's consistency is analysed in TSM-10 as logical consistency with the essential truths. Counting primitives therefore gives no advantage to the realist.
\end{reply}

\begin{objection}[Actuality]
Only modal realism explains why we are in the actual world, rather than merely in a world that represents itself as actual: \enquote{actual} is indexical, and the question has a trivial answer.
\end{objection}

\begin{reply}
On the ersatz account the question has a still more trivial answer. The actual world is the set of truths, and the cosmos is the only concrete totality, so there is nowhere else to be (TSM-07). Lewis's indexical account is needed only because there are other concrete totalities whose inhabitants call their worlds \enquote{actual} with equal right; it removes a problem that only modal realism creates.
\end{reply}

\begin{objection}[The analogy with sets]
We believe in sets because of their theoretical benefits, though they too are isolated from us. Parity requires the same for worlds.
\end{objection}

\begin{reply}
Parity requires the same test, and the two ontologies fare differently under it. No rival reconstructs classical mathematics at lower cost than set theory; the ersatz account does all the work of concrete worlds, using abstract objects that, like sets, we already accept.
\end{reply}

\section{Conclusion}

Modal realism asks us to believe in concrete universes isolated from us, in return for theoretical benefits. The analyses of counterfactuals, content and properties are available without them; the reduction of modality to the non-modal is not achieved with them. Other possible worlds do not exist concretely.

\begin{thebibliography}{99}\small
\bibitem[Divers and Melia(2002)]{DiversMelia2002} Divers, J., and Melia, J. (2002). The analytic limit of genuine modal realism. \emph{Mind}, 111, 15--36.
\bibitem[Lewis(1973)]{Lewis1973} Lewis, D. (1973). \emph{Counterfactuals}. Blackwell.
\bibitem[Lewis(1986)]{Lewis1986} Lewis, D. (1986). \emph{On the Plurality of Worlds}. Blackwell.
\end{thebibliography}

\end{document}
